\SourceLeaf{323}den gesuchten mit 0. Ferner bezeichne ich die Neigung der Geraden 12 zu einer festen Geraden $L$ einfach mit 12. Man betrachtet diese Neigung als in einem \textls[120]{bestimmten} Sinn von 0 bis $360^0$ zunehmend (zur Veranschaulichung denken Sie sich $L$ als Meridian und 12 als das Azimut von 2, von 1 aus gesehen). Entsprechend und auf dieselbe Gerade $L$ bezogen sind 21, 13, 31, 23, 32 zu verstehen. Dabei versteht sich von selbst, dass $21-12=\pm180^0$ (oder allgemeiner gleich einem ungeraden Vielfachen von $180^0$) ist. Die Zeichen 01, 02, 03 sollen eine ähnliche Bedeutung haben, mit dem Unterschied, dass sie sich statt auf $L$ auf \textls[120]{eine andere ganz willkürliche} feste Gerade beziehen.

Kurz könnte man sagen, dass 12, 13, 21, 23, 31, 32 die wahren Azimute der betreffenden Richtungen sind, 01, 02, 03 dagegen die noch nicht orientierten Azimute oder diejenigen, die einer noch unbekannten Orientierungskorrektur bedürfen. Azimut und Meridian dienen hier selbstverständlich nur der Veranschaulichung.

Das Kriterium dafür, ob die vorgeschriebenen 6 Größen (wenn Sie 12 und 21 nur als eine zählen) physisch verträglich sind, lautet wie folgt.
\[
\left.
\begin{array}{l}
\sin(12-13-02+03)\\
\sin(23-21-03+01)\\
\sin(31-32-01+02)
\end{array}
\right\}\ \text{müssen dasselbe Vorzeichen haben.}
\]

Das vorhin erwähnte Theorem über Vierecke, das meines Wissens\HistoricalNote{*)}{*)}{Sollte ich mich darin irren und sollten Sie mir nachweisen können, dass das Theorem schon jemand anderes aufgestellt hat, wäre ich Ihnen sehr verbunden.} bisher unbekannt ist und sich seit fast einem halben Jahrhundert in meinem Besitz befindet, lautet wie folgt.
\[
\frac{\sin(12-13-02+03)}{(01).(23)}=\frac{\sin(23-21-03+01)}{(02).(13)}=\frac{\sin(31-32-01+02)}{(03).(12)}\,,
\]

wobei die Faktoren in den Nennern die durchweg positiv genommenen Abstände zwischen den betreffenden Punkten ausdrücken. Die \textls[120]{Ableitung} dieses schönen Theorems mithilfe komplexer Zahlen ist denkbar einfach. Hier will ich sie nur andeuten.



